\chapter{Introduction}
The internship addressed a practical question in TSP production: whether historical, one-minute process measurements can support useful estimates or short-horizon forecasts of \texttt{SLURRY\_FREE\_ACID}. The work was developed as an industrial machine-learning project rather than as a single offline model. It therefore included data investigation, leakage-controlled experiments, model selection, an interactive dashboard, and a laboratory deployment path built around OPC UA and Docker.

The project has two motivations. First, a quality-related measurement may not be continuously available as an online signal. Second, even when a recent laboratory result is available, anticipating its movement can provide an earlier indication of a change. These are different operational problems with different allowable inputs. Treating them as one task would overstate what a model can do.

The general objective was to create a transparent, reproducible advisory prototype for free-acid monitoring. The specific objectives were to understand the available historian extract, prepare time-respecting features, compare appropriate baselines and models, distinguish future forecasting from current-time virtual sensing, test robustness across chronological folds, and connect saved artifacts to a safe read-only streaming architecture.

The main contributions are: a documented feature contract with lags, rolling descriptors and process proxies; a target-anchored delta forecast that outperformed persistence on the held-out January period; a separate target-free virtual-sensor line of work; stability-oriented clustered Ridge experiments; and an advisory-only CSV-replay implementation with quality checks, warm-up, persistence and monitoring.

The report follows the engineering progression of the project. Chapter 2 provides the available company context. Chapters 3 to 7 cover the problem, data, methodology, results and interpretation. Chapters 8 to 10 describe deployment, implementation and limitations. The final chapters state the contribution, conclusions and realistic next work.
